\documentclass[../main.tex]{subfiles}
\ifSubfilesClassLoaded{\setcounter{chapter}{6}}{}
\begin{document}

\chapter{Struktur Data dan Pemrosesan String Lanjut}

\begin{subcpmk}
  \item Sub-CPMK 4.1: Mengoperasikan \texttt{list}, \texttt{tuple}, \texttt{dict}, dan \texttt{set} secara lanjut (mutasi, metode, komposisi)
  \item Sub-CPMK 4.2: Memanipulasi string dengan \textit{slicing}, metode lanjut, dan pemformatan keluaran
\end{subcpmk}

\SesiRPS{8}{4.1 dan 4.2}{$\sim$170 menit}{CPMK-4}

\section{Sarana dan prasarana}

Python 3.10+, editor, terminal.

\section{Prosedur kerja}

\begin{enumerate}[leftmargin=*]
  \item Eksplorasi metode \texttt{list}, \texttt{dict}, \texttt{set}, \texttt{tuple} dari dokumentasi ringkas modul.
  \item Latih teknik iterasi lanjut dan manipulasi string.
  \item Selesaikan aktivitas dengan data terstruktur sederhana.
\end{enumerate}

\section{List: Lebih Dalam}

\texttt{list} adalah struktur data \textit{mutable} paling sering digunakan. Indeks dimulai dari 0; irisan \texttt{s[a:b]} eksklusif di ujung kanan.

\begin{lstlisting}[caption=Metode list lengkap (Tutorial Python ss5.1)]
angka = [3, 1, 4, 1, 5, 9, 2, 6]
angka.append(7)          # tambah di akhir
angka.insert(2, 99)      # sisipkan di indeks 2
angka.remove(1)          # hapus kemunculan pertama nilai 1
dikeluarkan = angka.pop()    # hapus & kembalikan elemen terakhir
dikeluarkan2 = angka.pop(0)  # hapus & kembalikan di indeks 0
angka.sort()             # urut ascending (in-place)
angka.reverse()          # balik urutan (in-place)
salinan = angka.copy()   # buat salinan dangkal
print(angka.count(1))    # hitung kemunculan nilai 1
print(angka.index(5))    # indeks pertama nilai 5
\end{lstlisting}

\subsection{List sebagai Stack (LIFO)}

Stack menggunakan \texttt{append()} untuk push dan \texttt{pop()} untuk pop (Tutorial Python §5.1.1):

\begin{lstlisting}[caption=List sebagai stack]
stack = []
stack.append("halaman_1")
stack.append("halaman_2")
stack.append("halaman_3")
print(stack.pop())   # "halaman_3" -- LIFO
print(stack.pop())   # "halaman_2"
\end{lstlisting}

\subsection{List sebagai Queue (FIFO)}

Untuk queue efisien, gunakan \texttt{collections.deque} karena \texttt{list.pop(0)} lambat untuk list besar (Tutorial Python §5.1.2):

\begin{lstlisting}[caption=Queue dengan collections.deque]
from collections import deque
antrian = deque(["pelanggan_1", "pelanggan_2", "pelanggan_3"])
antrian.append("pelanggan_4")     # masuk di belakang
print(antrian.popleft())          # keluar dari depan: "pelanggan_1"
print(antrian.popleft())          # "pelanggan_2"
\end{lstlisting}

\subsection{List Comprehension}

List comprehension menyediakan cara ringkas untuk membuat list (Tutorial Python §5.1.3):

\begin{lstlisting}[caption=List comprehension]
# Cara lama: loop biasa
kuadrat = []
for i in range(10):
    kuadrat.append(i ** 2)

# List comprehension -- lebih ringkas
kuadrat = [i ** 2 for i in range(10)]
print(kuadrat)   # [0, 1, 4, 9, 16, 25, 36, 49, 64, 81]

# Dengan kondisi filter
genap_kuadrat = [i ** 2 for i in range(10) if i % 2 == 0]
print(genap_kuadrat)   # [0, 4, 16, 36, 64]

# Manipulasi string
kata = ["  apel  ", "  JERUK  ", " mangga "]
bersih = [k.strip().lower() for k in kata]
print(bersih)   # ['apel', 'jeruk', 'mangga']
\end{lstlisting}

\section{Pernyataan \texttt{del}}

\texttt{del} menghapus elemen dari list berdasarkan indeks atau slice, atau menghapus variabel (Tutorial Python §5.2):

\begin{lstlisting}[caption=Pernyataan del]
a = [1, 2, 3, 4, 5, 6]
del a[0]       # hapus elemen pertama -> [2, 3, 4, 5, 6]
del a[2:4]     # hapus slice -> [2, 3, 6]
print(a)

# Menghapus variabel dari namespace
x = 100
del x
# print(x)  # akan menghasilkan NameError
\end{lstlisting}

\section{Tuple dan Sequences}

\texttt{tuple} tidak dapat diubah (\textit{immutable}). Berguna untuk data yang tidak boleh dimodifikasi (Tutorial Python §5.3):

\begin{lstlisting}[caption=Tuple: packing dan unpacking]
# Tuple packing
titik = 3, 7      # atau titik = (3, 7)

# Tuple unpacking
x, y = titik
print(f"x={x}, y={y}")

# Multiple assignment
a, b, c = 1, 2, 3

# Swap nilai tanpa variabel sementara
a, b = b, a
print(a, b)   # 2 1

# Tuple satu elemen memerlukan koma
satu = (42,)
print(type(satu))   # <class 'tuple'>
\end{lstlisting}

\section{Set}

Koleksi unik tanpa urutan tetap; berguna untuk menghapus duplikasi atau uji keanggotaan cepat (Tutorial Python §5.4):

\begin{lstlisting}[caption=Operasi set]
A = {1, 2, 3, 4, 5}
B = {4, 5, 6, 7, 8}

print(A | B)   # union:        {1,2,3,4,5,6,7,8}
print(A & B)   # intersection: {4, 5}
print(A - B)   # difference:   {1, 2, 3}
print(A ^ B)   # symmetric diff: {1,2,3,6,7,8}

# Cek keanggotaan O(1)
print(3 in A)     # True
print(9 not in A) # True

# Hapus duplikasi dari list
data = [1, 2, 2, 3, 3, 3, 4]
unik = list(set(data))
print(sorted(unik))   # [1, 2, 3, 4]
\end{lstlisting}

\section{Dictionary}

Memetakan kunci ke nilai; kunci harus \textit{hashable} (Tutorial Python §5.5):

\begin{lstlisting}[caption=Dictionary: operasi lanjut]
# Membuat dict
mahasiswa = {"nim": "12345", "nama": "Aulia", "ipk": 3.5}

# Akses aman dengan get (tidak memunculkan Error)
print(mahasiswa.get("alamat", "Tidak ada"))  # Tidak ada

# Iterasi
for kunci, nilai in mahasiswa.items():
    print(f"  {kunci}: {nilai}")

# Dict comprehension
kuadrat_dict = {x: x**2 for x in range(6)}
print(kuadrat_dict)  # {0:0, 1:1, 2:4, 3:9, 4:16, 5:25}

# Menggabungkan dua dict (Python 3.9+)
info1 = {"a": 1, "b": 2}
info2 = {"c": 3, "d": 4}
gabungan = info1 | info2
\end{lstlisting}

\section{Teknik Iterasi Lanjut}

Python menyediakan fungsi bawaan yang membuat iterasi lebih ekspresif (Tutorial Python §5.6):

\begin{lstlisting}[caption=enumerate zip dan items]
# enumerate: indeks + nilai sekaligus
buah = ["apel", "jeruk", "mangga"]
for i, b in enumerate(buah, start=1):
    print(f"{i}. {b}")

# zip: iterasi dua list bersamaan
nama = ["Andi", "Budi", "Citra"]
nilai = [85, 72, 91]
for n, v in zip(nama, nilai):
    print(f"{n}: {v}")

# sorted: mengembalikan list terurut
angka = [3, 1, 4, 1, 5, 9, 2]
print(sorted(angka))              # ascending
print(sorted(angka, reverse=True))# descending

# reversed: iterasi terbalik
for x in reversed(range(1, 6)):
    print(x, end=" ")   # 5 4 3 2 1
\end{lstlisting}

\section{String: Slicing dan Metode Lanjut}

String adalah sequence immutable. \texttt{f-string} memudahkan interpolasi dengan format:

\begin{lstlisting}[caption=f-string dengan format specification]
pi = 3.14159
nama = "Bayu"
print(f"Halo {nama}!")
print(f"Pi = {pi:.2f}")         # 2 desimal: 3.14
print(f"Pi = {pi:8.3f}")        # lebar 8, 3 desimal
print(f"Biner 42 = {42:b}")     # representasi biner
print(f"{'kiri':<10}|{'kanan':>10}")  # perataan
\end{lstlisting}

\section{Studi Kasus dan Contoh Kode}

Berikut adalah contoh-contoh program Python yang mengimplementasikan konsep-konsep pada modul ini.

\subsection{Kode: \texttt{matriks\_3x3.py}}
\begin{lstlisting}[language=Python,caption={\texttt{matriks\_3x3.py}}]
"""Matriks 3x3: isi dan cetak."""

# mat[i][j] untuk i,j in 1..3 - baris/kolom 0 tidak dipakai.


def isi(matrik: list[list[int]]) -> None:
    for i in range(1, 4):
        for j in range(1, 4):
            matrik[i][j] = int(input(f"A[{i},{j}] = "))


def cetak(matrik: list[list[int]]) -> None:
    for i in range(1, 4):
        for j in range(1, 4):
            print(f"{matrik[i][j]:3d}", end="")
        print()


def main() -> None:
    matrik = [[0] * 4 for _ in range(4)]
    print("Mengisi elemen matrik A")
    isi(matrik)
    print("Mencetak elemen matrik A")
    cetak(matrik)


if __name__ == "__main__":
    main()
\end{lstlisting}

\subsection{Kode: \texttt{array\_isi\_cetak.py}}
\begin{lstlisting}[language=Python,caption={\texttt{array\_isi\_cetak.py}}]
"""Array 1..3: isi dengan while, cetak dengan for."""

# Indeks 1..3 dipetakan ke list panjang 4 (slot 0 tidak dipakai).


def main() -> None:
    nilai = [0] * 4
    print("Mengisi elemen array")
    i = 1
    while i <= 3:
        nilai[i] = int(input(f"Nilai ke {i} = "))
        i += 1
    print("Mencetak elemen array")
    for i in range(1, 4):
        print(f"Nilai ke {i} = ", nilai[i])


if __name__ == "__main__":
    main()
\end{lstlisting}

\subsection{Kode: \texttt{array\_prosedur.py}}
\begin{lstlisting}[language=Python,caption={\texttt{array\_prosedur.py}}]
"""Prosedur isi dan cetak array [1..3]."""


def isi(nilai: list[int]) -> None:
    i = 1
    while i <= 3:
        nilai[i] = int(input(f"Nilai ke {i} = "))
        i += 1


def cetak(nilai: list[int]) -> None:
    for i in range(1, 4):
        print(f"Nilai ke {i} = ", nilai[i])


def main() -> None:
    nilai = [0] * 4
    print("Mengisi elemen array")
    isi(nilai)
    print("Mencetak elemen array")
    cetak(nilai)


if __name__ == "__main__":
    main()
\end{lstlisting}

\subsection{Kode: \texttt{daf\_pegawai\_array.py}}
\begin{lstlisting}[language=Python,caption={\texttt{daf\_pegawai\_array.py}}]
"""Daftar 5 pegawai: input lalu cetak satu per satu."""

from dataclasses import dataclass


@dataclass
class Pegawai:
    nama: str = ""
    gol: str = " "
    jml_anak: int = 0


def input_pegawai(pegawai: Pegawai) -> None:
    pegawai.nama = input("Nama        : ")
    pegawai.gol = input("Gol (A - D) : ").strip().upper()[:1]
    pegawai.jml_anak = int(input("Jumlah Anak : "))


def upah_kotor(gol: str) -> float:
    g = gol.upper()
    if g == "A":
        return 1_000_000.0
    if g == "B":
        return 800_000.0
    if g == "C":
        return 600_000.0
    if g == "D":
        return 400_000.0
    return 0.0


def persen_tunjangan(jml_anak: int) -> float:
    if jml_anak > 2:
        return 0.3
    return 0.2


def cetak(pegawai: Pegawai) -> None:
    print("Nama        :", pegawai.nama)
    print("Gol (A - D) :", pegawai.gol)
    print("Jumlah Anak :", pegawai.jml_anak)
    uk = upah_kotor(pegawai.gol)
    upah_bersih = uk - (uk * persen_tunjangan(pegawai.jml_anak))
    print(f"Upah        : {upah_bersih:5.2f}")


def main() -> None:
    daf_pegawai = [Pegawai() for _ in range(5)]
    for i in range(5):
        print("Data ke ", i + 1, ":", sep="")
        input_pegawai(daf_pegawai[i])
    for i in range(5):
        print("Data ke ", i + 1, ":", sep="")
        cetak(daf_pegawai[i])
        input("Tekan Enter untuk melanjutkan...")


if __name__ == "__main__":
    main()
\end{lstlisting}

\subsection{Kode: \texttt{record\_upah.py}}
\begin{lstlisting}[language=Python,caption={\texttt{record\_upah.py}}]
"""Record pegawai: upah kotor menurut gol dan tunjangan anak."""

from dataclasses import dataclass


@dataclass
class Pegawai:
    nama: str = ""
    gol: str = " "
    jml_anak: int = 0


def main() -> None:
    pegawai = Pegawai()
    pegawai.nama = input("Nama        : ")
    pegawai.gol = input("Gol (A - D) : ").strip().upper()[:1]
    pegawai.jml_anak = int(input("Jumlah Anak : "))

    if pegawai.gol == "A":
        upah_kotor = 1_000_000.0
    elif pegawai.gol == "B":
        upah_kotor = 800_000.0
    elif pegawai.gol == "C":
        upah_kotor = 600_000.0
    elif pegawai.gol == "D":
        upah_kotor = 400_000.0
    else:
        upah_kotor = 0.0

    if pegawai.jml_anak > 2:
        persen_tunjangan = 0.3
    else:
        persen_tunjangan = 0.2

        upah_bersih = upah_kotor - (upah_kotor * persen_tunjangan)
    print("Upah        :", upah_bersih)


if __name__ == "__main__":
    main()
\end{lstlisting}

\subsection{Kode: \texttt{record\_prosedur.py}}
\begin{lstlisting}[language=Python,caption={\texttt{record\_prosedur.py}}]
"""Input dan cetak pegawai dengan fungsi upah (setara record_prosedur.pas)."""

from dataclasses import dataclass


@dataclass
class Pegawai:
    nama: str = ""
    gol: str = " "
    jml_anak: int = 0


def input_pegawai(pegawai: Pegawai) -> None:
    pegawai.nama = input("Nama        : ")
    pegawai.gol = input("Gol (A - D) : ").strip().upper()[:1]
    pegawai.jml_anak = int(input("Jumlah Anak : "))


def upah_kotor(gol: str) -> float:
    g = gol.upper()
    if g == "A":
        return 1_000_000.0
    if g == "B":
        return 800_000.0
    if g == "C":
        return 600_000.0
    if g == "D":
        return 400_000.0
    return 0.0


def persen_tunjangan(jml_anak: int) -> float:
    # Pascal memakai nama variabel campur huruf pada badan fungsi
    if jml_anak > 2:
        return 0.3
    return 0.2


def cetak(pegawai: Pegawai) -> None:
    print("Nama        :", pegawai.nama)
    print("Gol (A - D) :", pegawai.gol)
    print("Jumlah Anak :", pegawai.jml_anak)
    uk = upah_kotor(pegawai.gol)
    upah_bersih = uk - (uk * persen_tunjangan(pegawai.jml_anak))
    print(f"Upah        : {upah_bersih:5.2f}")


def main() -> None:
    pegawai = Pegawai()
    input_pegawai(pegawai)
    cetak(pegawai)


if __name__ == "__main__":
    main()
\end{lstlisting}

\section{Referensi sesi}

\begin{itemize}[leftmargin=*]
  \item \textbf{[15]} Google, \textit{Python Class}.
  \item \textbf{[8, 9]} \textit{The Python Tutorial}, struktur data.
\end{itemize}

\begin{aktivitas}
  \item Diberikan list nilai, hitung rata-rata menggunakan list comprehension untuk menyaring nilai di atas rata-rata.
  \item Hitung frekuensi kata sederhana pada string menggunakan \texttt{dict} dan \texttt{split()}.
  \item Konversi list yang berisi duplikasi menjadi \texttt{set}, lalu kembali ke list terurut.
  \item Gunakan \texttt{enumerate} untuk mencetak menu bernomor dari list pilihan.
  \item Buat dict comprehension yang memetakan nama buah ke panjang namanya.
  \item Implementasikan stack sederhana dengan list untuk mensimulasikan tombol ``undo'' editor teks.
\end{aktivitas}

\begin{latihan}
  \item Jelaskan perbedaan \texttt{list} dan \texttt{tuple}. Kapan \texttt{tuple} lebih tepat?
  \item Apa yang terjadi jika kunci dict tidak ada saat diakses dengan \texttt{[]}? Bagaimana cara aman mengaksesnya?
  \item Tulis list comprehension yang menghasilkan angka ganjil dari 1 sampai 20.
  \item Apa perbedaan antara \texttt{list.sort()} dan \texttt{sorted(list)}?
  \item Refleksi: kapan \texttt{set} lebih tepat daripada \texttt{list}?
\end{latihan}

\begin{asesmen}
\textbf{Soal singkat}

\begin{enumerate}
  \item List comprehension \texttt{[x*2 for x in range(5) if x\%2==0]} menghasilkan:
  \begin{enumerate}
    \item \texttt{[0, 4, 8]}
    \item \texttt{[0, 2, 4, 6, 8]}
    \item \texttt{[2, 4, 6, 8, 10]}
    \item \texttt{[0, 1, 4]}
  \end{enumerate}
  \item Ekspresi manakah yang menghasilkan \texttt{[2, 4]}?
  \begin{enumerate}
    \item \texttt{[1,2,3,4][1:3]}
    \item \texttt{[1,2,3,4][1:4]}
    \item \texttt{[1,2,3,4][0:2]}
    \item \texttt{[1,2,3,4][2]}
  \end{enumerate}
\end{enumerate}

\textbf{Tugas}: program yang membaca beberapa baris (NIM,nama,nilai), menyimpan ke \texttt{dict} dengan kunci NIM, dan mencetak rekap termasuk rata-rata nilai kelas. Gunakan list comprehension untuk mendapatkan daftar mahasiswa dengan nilai di atas rata-rata.
\end{asesmen}

\begin{checklist}
  \item Saya dapat membuat dan mengubah \texttt{list} dengan metode seperti \texttt{append}, \texttt{pop}, \texttt{sort}
  \item Saya dapat menggunakan list comprehension
  \item Saya memahami \texttt{del} untuk menghapus elemen dan variabel
  \item Saya dapat mensimulasikan stack dan queue dengan \texttt{list} dan \texttt{deque}
  \item Saya dapat memakai \texttt{dict} untuk struktur data asosiatif
  \item Saya dapat menggunakan \texttt{set} untuk operasi himpunan dan menghapus duplikasi
  \item Saya dapat memakai \texttt{enumerate} dan \texttt{zip} dalam iterasi
  \item Saya dapat memformat string dengan \texttt{f-string} dan format specification
\end{checklist}

\begin{rangkuman}
Struktur data memilih representasi yang tepat untuk masalah: urutan mutable (\texttt{list}), urutan tidak berubah (\texttt{tuple}), pemetaan kunci-nilai (\texttt{dict}), koleksi unik (\texttt{set}). List comprehension dan fungsi \texttt{enumerate}/\texttt{zip} membuat iterasi lebih ringkas dan ekspresif.

\textbf{Referensi:} {\small Python Tutorial §5 \url{https://docs.python.org/3/tutorial/datastructures.html}}

\textbf{Kata kunci:} \code{list}, \code{tuple}, \code{dict}, \code{set}, \code{del}, \code{enumerate}, \code{zip}, \code{sorted}, \code{f-string}, list comprehension
\end{rangkuman}

\end{document}
